\section{Canonical reduced standard bases of the actual formal ideal}

Let $S=\mathbb C[[X_1,\ldots,X_n]]$, let $\prec$ be a monomial
order, and let $\mathcal N(I)$ be the actual initial-exponent diagram
of an arbitrary ideal $I\subset S$. Write $R_I$ for the remainder
already constructed by actual formal division. Its coefficients vanish
on $\mathcal N(I)$. The vertex members chosen for that division are
monic but need not be reduced. We now construct reduced members from the
remainder itself. This is the normal-form construction used in the
standard-basis argument in \cite{bierstone-milman03-desingularization-algorithms},
Section 6.3; here existence, uniqueness and generation are established
directly for the actual formal ideal, before any geometric hypothesis.

\begin{theorem}[The actual remainder is the unique normal representative]
\label{hir:actual-formal-normal-representative-unique}
\uses{hir:actual-formal-vertex-division-identity,hir:actual-formal-vertex-division-supports}
\source{bierstone-milman03-desingularization-algorithms}
\lean{Hironaka.HironakaActualReducedFormalStandardBasis.sub_actualFormalRemainder_mem,Hironaka.HironakaActualReducedFormalStandardBasis.eq_zero_of_mem_of_diagram_coefficients_zero,Hironaka.HironakaActualReducedFormalStandardBasis.actualFormalRemainder_unique,Hironaka.HironakaActualReducedFormalStandardBasis.actualFormalRemainder_eq_zero_iff,Hironaka.HironakaActualReducedFormalStandardBasis.actualFormalRemainder_add,Hironaka.HironakaActualReducedFormalStandardBasis.actualFormalRemainder_smul,Hironaka.HironakaActualReducedFormalStandardBasis.actualFormalRemainder_eq_self,Hironaka.HironakaActualReducedFormalStandardBasis.actualFormalRemainder_idempotent,Hironaka.HironakaActualReducedFormalStandardBasis.actualFormalNormalProjection,Hironaka.HironakaActualReducedFormalStandardBasis.mem_range_actualFormalNormalProjection_iff}
\leanok
For every $F\in S$, $F-R_I(F)\in I$. If $F-H\in I$ and every
coefficient of $H$ on $\mathcal N(I)$ is zero, then $H=R_I(F)$.
Consequently $R_I(F)=0$ if and only if $F\in I$, and $R_I$ is a
$\mathbb C$-linear idempotent with kernel $I$ and image the space of
series supported outside $\mathcal N(I)$.
\end{theorem}
\begin{proof}
The division identity expresses $F-R_I(F)$ as a finite sum of multiples
of actual members of $I$, proving membership. The difference
$H-R_I(F)=(F-R_I(F))-(F-H)$ belongs to $I$ and has zero coefficients
on its entire diagram. If it were nonzero, its actual initial exponent
would belong to $\mathcal N(I)$, contradicting that vanishing. This proves
uniqueness. A member of $I$ has zero remainder by the same argument;
conversely a zero remainder and the membership just proved put $F$ in $I$.

For $F,G\in S$ and $c\in\mathbb C$, the series $R_I(F)+R_I(G)$ and
$cR_I(F)$ vanish on the diagram and differ from $F+G$ and $cF$ by members
of $I$. Uniqueness proves additivity and homogeneity. Any series vanishing
on the diagram is its own remainder, again by uniqueness, proving both
idempotence and the asserted image. No decomposition is supplied by the
caller.
\end{proof}

\begin{lemma}[Actual division preserves every initial segment]
\label{hir:actual-formal-division-initial-segment}
\uses{hir:actual-formal-vertex-division-identity}
\lean{Hironaka.HironakaActualReducedFormalStandardBasis.actualDivisionCoefficient_eq_zero_of_lt,Hironaka.HironakaActualReducedFormalStandardBasis.actualFormalRemainder_coefficient_lower}
\leanok
If $[X^c]F=0$ for every $c\prec a$, then $[X^c]R_I(F)=0$ for every
$c\prec a$.
\end{lemma}
\begin{proof}
Use well-founded induction on $c$ in the actual division recursion.
For $c\prec a$, its input coefficient is zero. Every correction summand
contains a division coefficient at an exponent $b\prec c$, and therefore
$b\prec a$; the induction hypothesis makes that summand zero. The
recursion thus makes the division coefficient at $c$ zero. A remainder
coefficient is either this division coefficient or zero according as
$c$ lies outside or inside the actual diagram, proving the claim for all
initial segments, including monomial orders with infinitely many
predecessors.
\end{proof}

\begin{theorem}[Internally constructed canonical reduced vertex members]
\label{hir:actual-reduced-formal-vertex-members}
\uses{hir:actual-formal-normal-representative-unique,hir:actual-formal-division-initial-segment}
\lean{Hironaka.HironakaActualReducedFormalStandardBasis.actualReducedVertexSeries,Hironaka.HironakaActualReducedFormalStandardBasis.actualReducedVertexSeries_mem,Hironaka.HironakaActualReducedFormalStandardBasis.actualReducedVertexSeries_coefficient_diagram,Hironaka.HironakaActualReducedFormalStandardBasis.actualReducedVertexSeries_coefficient,Hironaka.HironakaActualReducedFormalStandardBasis.actualReducedVertexSeries_coefficient_lower,Hironaka.HironakaActualReducedFormalStandardBasis.actualReducedVertexSeries_ne_zero,Hironaka.HironakaActualReducedFormalStandardBasis.actualReducedVertexSeries_initialExponent,Hironaka.HironakaActualReducedFormalStandardBasis.actualReducedVertexSeries_unique}
\leanok
For each actual vertex $a$ of $\mathcal N(I)$, define
\[
 G_a=X^a-R_I(X^a).
\]
Then $G_a\in I$, its actual initial exponent is $a$, its coefficient
there is one, and its coefficient at every other exponent of
$\mathcal N(I)$ is zero. It is the unique member of $I$ with this entire
diagram coefficient vector. In particular it depends only on the actual
ideal and monomial order.
\end{theorem}
\begin{proof}
Membership follows from the first theorem. At a diagram exponent $b$,
the remainder coefficient is zero, so $[X^b]G_a$ is the coefficient of
$X^a$, namely one for $b=a$ and zero otherwise. All coefficients of
$X^a$ below $a$ vanish; the initial-segment lemma makes the remainder
coefficients there vanish as well. Thus the nonzero coefficient at $a$
is the actual first one. If another member has the same coefficients
on the whole diagram, its difference from $G_a$ belongs to $I$ and
vanishes on the diagram, so the first proof makes that difference zero.
This also proves independence from all choices made during division.
\end{proof}

\begin{theorem}[The genuine reduced standard basis generates the whole ideal]
\label{hir:actual-reduced-formal-standard-basis-generation}
\uses{hir:actual-reduced-formal-vertex-members,hir:actual-formal-normal-representative-unique,hir:actual-formal-ideal-initial-diagram-upper}
\lean{Hironaka.HironakaActualReducedFormalStandardBasis.ideal_eq_of_le_of_diagram_le,Hironaka.HironakaActualReducedFormalStandardBasis.actual_formal_ideal_eq_span_reduced_vertex_series}
\leanok
The finite family $(G_a)$ is a reduced standard basis of the SAME ideal:
\[
 I=(G_a:a\text{ an actual vertex of }\mathcal N(I)).
\]
This includes the zero ideal, the unit ideal and ambient dimension zero.
\end{theorem}
\begin{proof}
Let $J$ be the displayed span. Its generators belong to $I$, hence
$J\subset I$. Each actual vertex $a$ of $\mathcal N(I)$ is the actual
initial exponent of $G_a\in J$. Upward closure of the actual diagram and
the vertex cover imply $\mathcal N(I)\subset\mathcal N(J)$; the reverse
inclusion follows directly from $J\subset I$.

For any $F\in I$, division by the actual ideal $J$ gives
$R_J(F)=F-H$ with $H\in J\subset I$. Thus $R_J(F)\in I$.
If it were nonzero, its actual initial exponent would lie in
$\mathcal N(I)\subset\mathcal N(J)$, where all remainder coefficients
vanish, a contradiction. Therefore $R_J(F)=0$ and the first theorem
gives $F\in J$. This proves the other whole-ideal inclusion. No
geometric compatibility, normal-flatness, or strict-transform conclusion
has been assumed or asserted.
\end{proof}
